\section{How Evidence Debt Is Created (RQ1)}
\label{sec:conditions}

Evidence debt is not created by exotic failures. It is created by ordinary, often
locally rational practices at change time, and by environmental properties of
cloud-native platforms that make evidence perishable by default. We analyze both,
because the distinction matters for remediation: practices can be changed; platform
properties must be compensated for.

\subsection{The Evidence a Change Should Leave Behind}
\label{sec:schema-informal}

Fix, informally for now (formally in \cref{def:schema}), the six interrogatives a
complete operational record answers for a change: \emph{who} acted (a human or
machine principal, resolvable to an accountable identity); \emph{what} changed (the
content delta, identified precisely enough to reproduce it); \emph{why} (the intent:
the problem, requirement, or decision the change served); \emph{under which
authorization} (the approval or policy basis that legitimized it); \emph{with which
outcome} (whether it succeeded, what it affected, what followed); and \emph{how the
resulting state can be reconstructed} (the recovery-relevant linkage from running
state back to buildable sources). Contemporary toolchains capture fragments of these
in different systems---tickets, commit messages, review records, CI logs, registries,
deployment tools, cloud audit trails~\cite{cloudtrail,ghdocs}---keyed by correlation
identifiers whose continuity is nobody's explicit responsibility~\cite{sigelman2010,
otel}.

\subsection{Debt-Creating Practices}
\label{sec:practices}

The following practices, enumerated from the structure of the six interrogatives
(each practice deletes or degrades at least one), create deficits at change time.
We state them descriptively, not judgmentally; several are unavoidable under
incident pressure, which is precisely why the debt frame---deliberate deferral with
a later price---fits them.

\begin{description}[leftmargin=1.6em,itemsep=1.5pt,topsep=2pt]
\item[P1: Ticketless change.] Work performed without an intent record, or with intent
recorded only in ephemeral channels (chat, hallway); deletes \emph{why}.
\item[P2: Out-of-band execution.] Manual changes applied directly to runtime
(console edits, \texttt{kubectl} against production, emergency patches) that bypass
the pipeline where capture is instrumented; degrades \emph{what}, \emph{who}, and
\emph{outcome} simultaneously.
\item[P3: Shared and impersonal identity.] Changes executed under shared accounts,
service credentials used interactively, or broad break-glass roles; deletes
\emph{who} even when an API trail exists~\cite{cloudtrail}.
\item[P4: Approval as gesture.] Authorization expressed as transient platform state
(a review click, a pipeline gate) not bound durably to the executed content; deletes
\emph{under which authorization} on governance timescales.
\item[P5: Correlation rupture.] Failure to propagate identifiers across stages---no
ticket reference in the commit, no commit in the build record, no digest in the
deployment record; leaves all facts individually present but \emph{unjoinable},
which \cref{sec:simulation} shows is a distinct and subtler deficit than absence.
\item[P6: Retention minimization by default.] Accepting platform default retention
(often 30--90 days for logs~\cite{nist80092,ghdocs,cloudtrail}) for records whose
governance-relevant lifetime is years; converts captured evidence into future
deficits on a schedule.
\item[P7: Migration without evidence portability.] Tool and vendor changes (forge,
ITSM, CI) that migrate active data but orphan or discard historical records and
their identifiers; a bulk deficit-creation event.
\item[P8: Knowledge kept in people.] Reliance on the memory of specific engineers as
the de facto index into operational history; converts turnover~\cite{rigby2016,
avelino2016} into evidence decay.
\end{description}

\subsection{Environmental Amplifiers}
\label{sec:amplifiers}

Cloud-native platforms amplify the consequences of P1--P8 in ways traditional
infrastructure did not. \emph{Ephemerality}: containers, nodes, and serverless
instances are routinely destroyed and replaced~\cite{burns2016}; state that would
once have persisted incidentally on long-lived servers (shell histories, local logs,
file timestamps) now vanishes by design, removing the accidental evidence that
forensic practice historically relied on~\cite{casey2011}. \emph{Volume and
automation}: reconciliation loops and autoscalers generate changes at rates no
manual record-keeping can track, so anything not captured automatically is not
captured~\cite{beetz2022,opengitops}. \emph{Fragmentation}: the record of one
logical change is scattered across six or more stores with heterogeneous schemas,
identifiers, clocks, and retention policies---the precondition for the correlation
failures that \cref{sec:simulation} quantifies. \emph{Schema and identity drift}:
service renames, account migrations, and observability re-instrumentation change the
vocabulary in which old evidence is expressed, so even surviving records become
progressively harder to join (a decay channel distinct from deletion,
\cref{sec:interest}).

\subsection{Why the Cost Is Deferred, Not Avoided}
\label{sec:deferral}

Each practice above trades a certain, small, present cost (filing the ticket,
propagating the identifier, exporting before migration) against an uncertain, larger,
future cost that is realized only if a \emph{query} against operational history
arrives: an incident, an audit, an investigation, a recovery, a dispute. This
contingent structure is what makes the debt frame apt and what makes the phenomenon
invisible to steady-state operational metrics: a system can run flawlessly for years
while its reconstructable history rots. Whether the deferred cost is eventually paid
is therefore a question about the query workload (how often history is interrogated,
with what stakes) and about decay (what survives until then)---exactly the two
quantities the formal model of \cref{sec:model} is built around, and the reason the
central claim of \cref{sec:intro} is a testable hypothesis rather than a tautology:
if queries are rare and redundancy is high, deferral could in principle be nearly
free. \Cref{sec:simulation} tests when it is and is not.

\begin{figure}[t]
\centering
\resizebox{\columnwidth}{!}{% Evidence-debt lifecycle (single column, resized)
\begin{tikzpicture}[node distance=3.6mm and 4.5mm]
  % phase headers
  \node[lbl, anchor=west] (h1) at (0,0) {\textbf{Change time}};
  \node[lbl, anchor=west] (h2) at (3.45,0) {\textbf{Latency period}};
  \node[lbl, anchor=west] (h3) at (7.4,0) {\textbf{Query time}};

  % change time
  \node[human, below=3mm of h1.west, anchor=north west, minimum width=26mm] (act)
    {change occurs\\{\smlbl who / what / why / authz}\\[-1pt]{\smlbl outcome / recovery}};
  \node[attn, below=3mm of act, minimum width=26mm] (cap)
    {captured evidence\\{\smlbl records $+$ link keys}};
  \node[brok, below=3mm of cap, minimum width=26mm] (def)
    {\textbf{deficit} (principal)\\{\smlbl P1--P8: deferred capture}};
  \draw[flow] (act) -- (cap);
  \draw[bflow] (act.south east) to[out=-40,in=40] (def.north east);

  % latency
  \node[evid, right=6mm of cap, minimum width=27mm] (decay)
    {decay processes\\{\smlbl retention expiry $\cdot$ migration}\\[-1pt]{\smlbl turnover $\cdot$ ephemerality $\cdot$ id drift}};
  \node[brok, below=3mm of decay, minimum width=27mm] (interest)
    {\textbf{interest}\\{\smlbl recovery paths die;}\\[-1pt]{\smlbl cost plateaus $\to$ cliffs}};
  \draw[eflow] (cap) -- (decay);
  \draw[bflow] (def.east) to[out=0,in=200] (interest.west);
  \draw[flow] (decay) -- (interest);

  % query time
  \node[human, right=6mm of decay, minimum width=25mm] (query)
    {query arrives\\{\smlbl incident $\cdot$ audit $\cdot$ investigation}\\[-1pt]{\smlbl recovery $\cdot$ litigation}};
  \node[stage, below=3mm of query, minimum width=25mm] (rec)
    {reconstruction\\{\smlbl effort (ERC)}\\[-1pt]{\smlbl error (false joins)}};
  \node[brok, below=3mm of rec, minimum width=25mm] (wo)
    {\textbf{write-off}\\{\smlbl irrecoverable facts}\\[-1pt]{\smlbl realized liability}};
  \draw[flow] (query) -- (rec);
  \draw[bflow] (rec) -- (wo);
  \draw[eflow] (interest.east) to[out=0,in=190] (rec.west);
\end{tikzpicture}
}
\caption{The evidence-debt lifecycle. At change time, capture practices (top) either
record or defer each evidence element; deferral creates a deficit (principal). During
the latency period, decay processes---retention expiry, turnover, migration,
identifier drift---degrade the remaining reconstruction paths (interest). The debt is
realized when a query against operational history arrives (incident, audit,
investigation, recovery, litigation): surviving evidence is assembled, missing links
are reconstructed at effort cost, wrongly reconstructed at error cost, or found
irrecoverable (write-off). Boxes name the constructs formalized in
\cref{sec:definitions}.}
\label{fig:lifecycle}
\end{figure}

\Cref{fig:lifecycle} summarizes the lifecycle: creation at change time, silent
accumulation and decay during the latency period, and realization at query time.
The remainder of the paper formalizes each stage.
